% ==============================================================================
% Chapter 3: System Architecture & Agent Design
% Target Page Budget: ~10 pages
% ==============================================================================
\chapter{System Architecture \& Agent Design}
\label{chap:system_architecture}

\section{System Philosophy and Conceptual Blueprint}
\label{sec:arch_philosophy}
% [Drafting Guidance: ~1.5 pages]
% - Design principles: Non-judgmental reflection, user autonomy, cognitive transparency, privacy-by-design.
% - High-level conceptual flow: Ingestion -> Analysis -> Modeling -> Reflection -> Feedback.

\section{High-Level Architectural Topology}
\label{sec:arch_topology}
% [Drafting Guidance: ~2.0 pages]
% - Layered architecture diagram description (Data Layer, Cognitive Agent Layer, Presentation/Reflection Layer).
% - Event-driven communication bus and message brokerage across agents.
% - Integration of asynchronous worker queues and vector memory stores.

\section{Collaborative Multi-Agent Society: Roles and Responsibilities}
\label{sec:arch_agent_roles}
% [Drafting Guidance: ~3.0 pages]
% Detailed breakdown of the specialized agent roles:
% 1. \agentObserver{}: Stream parsing, modal feature extraction, metadata tagging.
% 2. \agentAnalyst{}: Semantic classification of comparison triggers, influencer disclosure detection, exposure intensity computation.
% 3. \agentReflector{}: Longitudinal cognitive state tracking, memory consolidation, user mirror generation.
% 4. \agentInterventionist{}: Adaptive nudge scheduling, cognitive reframing prompts, digital pause suggestions.

\subsection{The Observer Agent (\agentObserver)}
% [Detailed input/output specification, memory footprint, tools]

\subsection{The Analyst Agent (\agentAnalyst)}
% [Classification rubrics, prompt engineering templates, sentiment/comparison scoring]

\subsection{The Reflector Agent (\agentReflector)}
% [Cognitive synthesis, daily/weekly reflection generation, conversational interface]

\subsection{The Intervention Agent (\agentInterventionist)}
% [Behavioral economic nudges, reactance minimization, timing optimization]

\section{Multimodal Content Ingestion and Feature Extraction}
\label{sec:arch_ingestion_pipeline}
% [Drafting Guidance: ~1.5 pages]
% - Textual caption analysis and hashtag decomposition.
% - Visual feature heuristics (lifestyle aesthetics, luxury branding, body composition).
% - Engagement metadata (likes, comments, share velocity).

\section{Mathematical Modeling of Cognitive and Affective States}
\label{sec:arch_math_model}
% [Drafting Guidance: ~1.0 page]
% - Formal definition of the well-being state vector \wellbeingVector{}.
% - Equations defining social comparison salience \comparisonScore{} and influencer exposure index \exposureIndex{}.
% - Decay functions and temporal weighting in agent memory.

\section{Human-Agent Interaction and Explainability Interfaces}
\label{sec:arch_hai}
% [Drafting Guidance: ~1.0 page]
% - User dashboard wireframe concepts: Cognitive mirror view, exposure radar, reflective dialogue.
% - Explainable AI (XAI) principles in agent recommendations.
